\chapter{Компонентные граничные циклы и отсутствие связывания классов двенадцать и три}
\label{ch:v6-spectral-transition-component-boundary-gate}

\section{Уточнённый вопрос}

Предыдущий гейт установил физическое разложение
\begin{equation}
 H_{15}=H_q\oplus H_\ell,
 \qquad
 \dim_{\mathbb C}H_q=12,
 \qquad
 \dim_{\mathbb C}H_\ell=3.
 \label{eq:v6-component-boundary-split}
\end{equation}
Из аддитивности ранга ещё не следует, что оба слагаемых являются
самостоятельными спектральными циклами. Кроме того, даже при положительном
ответе полный класс пятнадцать может либо динамически связывать их в одном
ядре, либо фиксировать только суммарный заряд без совместной локализации.

Настоящий гейт различает четыре утверждения:
\begin{enumerate}
 \item формальное разложение целого числа $15=12+3$;
 \item существование отдельных комплексных и Real-граничных циклов;
 \item необходимость их суммы для исходного глобального класса;
 \item наличие энергии, связывающей две компоненты в одну частицу.
\end{enumerate}

\section{Редуцирующие конечные подциклы}

Пусть $P_q$ и $P_\ell$ --- проекторы из предыдущего гейта. Для текущего
заряженного пакета они удовлетворяют
\begin{equation}
 [P_s,a]=[P_s,D_{\mathrm{ch}}]=[P_s,\gamma]=0,
 \qquad s\in\{q,\ell\},
\end{equation}
а после ориентационного удвоения проектор
\begin{equation}
 P_s^{\mathbb R}=P_s\oplus\overline{P_s}
\end{equation}
коммутирует с обменной Real-структурой $J$.

Поэтому условие первого порядка наследуется сжатием:
\begin{equation}
 P_s\bigl[[D,a],JbJ^{-1}\bigr]P_s=0.
\end{equation}
Иными словами, кварковая и лептонная компоненты являются настоящими
редуцирующими подциклами замороженного заряженного $H_{15}$, а не только
разбиением списка состояний. Это соответствует диаграммному чтению
конечных спектральных троек, где компоненты графа задают редуцирующие
представления \cite{component-boundary-krajewski1998}.

Ограничение относится именно к текущему оператору с рёбрами $u,d,e$.
Невыведенный нейтринный оператор высшей степени сюда не добавляется.

\section{Явные тёплицевы символы компонент}

В коэффициентном пространстве $\mathbb C^{105}$ имеем ортогональные
проекторы
\begin{equation}
 q_0=q_q+q_\ell,
 \qquad
 q_qq_\ell=0,
 \qquad
 \rank q_q=12,
 \qquad
 \rank q_\ell=3.
\end{equation}
Для каждой компоненты определим Real-парный символ
\begin{align}
 V_s^+(z)&=zq_s+(1-q_s),\\
 V_s^-(z)&=z^{-1}\overline{q_s}+(1-\overline{q_s}).
 \label{eq:v6-component-boundary-ko7-symbol}
\end{align}
Обе матрицы унитарны, а обменная инволюция переводит $V_s^+$ в $V_s^-$.
Их комплексные граничные индексы равны
\begin{align}
 (\operatorname{ind}T_{V_q^+},\operatorname{ind}T_{V_q^-})&=(-12,+12),\\
 (\operatorname{ind}T_{V_\ell^+},\operatorname{ind}T_{V_\ell^-})&=(-3,+3).
\end{align}
В обменной вещественной форме соответствующие классы равны
\begin{equation}
 \kappa_q=12,
 \qquad
 \kappa_\ell=3
 \quad\text{в}\quad
 KO_6(M_{105}(\mathbb C)_{\mathbb R})\simeq\mathbb Z.
\end{equation}
Такое Real-чтение опирается на $KR$-теорию Атьи
\cite{component-boundary-atiyah1966}. Аддитивность граничного отображения
следует из тёплицевой точной последовательности
\cite{component-boundary-pimsner-voiculescu1980}.

\section{Неограниченные компонентные циклы}

Оператор числа и сдвиг ограничиваются без изменения формул:
\begin{equation}
 \mathcal H_s=\ell^2(\mathbb Z)\otimes q_s\mathbb C^{105},
 \qquad
 N_se_n=ne_n,
 \qquad
 U_se_n=e_{n+1},
 \qquad
 [N_s,U_s]=U_s.
\end{equation}
Резольвента $N_s$ компактна с конечной кратностью $\rank q_s$, а
коммутатор ограничен. Поэтому обе компоненты дают настоящие
неограниченные тёплицевы циклы в смысле неограниченного представления
$KK$-теории \cite{component-boundary-baaj-julg1983}.

Машинный аудит конечного окна дал нулевой остаток $[N_s,U_s]-U_s$ для
обоих рангов. Максимальные остатки унитарности символов равны
$3.85\cdot10^{-16}$ и $1.93\cdot10^{-16}$, а Real-обмена --- точному
нулю.

\section{Почему отдельный цикл не заменяет исходный дефект}

Положительный результат для подциклов не отменяет глобальный ledger:
\begin{equation}
 \kappa_{15}=15=\kappa_q+\kappa_\ell.
 \label{eq:v6-component-boundary-global-sum}
\end{equation}
Отдельный класс двенадцать или три не равен исходному классу пятнадцать и
не насыщает исходную нижнюю границу дефекта. Только прямая сумма
восстанавливает одновременно
\begin{align}
 -15&=-12-3,\\
 +15&=+12+3,\\
 \frac17&=\frac4{35}+\frac1{35}.
\end{align}

Следовательно, физическая редуцируемость и глобальная обязательность не
противоречат друг другу. Полный класс может быть составным, но его общая
величина всё равно фиксируется исходным замыканием.

\section{Проверка связывания}

Внутренний петлевой функционал имеет следовой вид
\begin{equation}
 S_{\mathrm{loop}}(V)=\tau_{210}([N,V]^*[N,V]).
\end{equation}
Для ортогональных $q_q,q_\ell$ смешанный след точно равен нулю. Поэтому
\begin{equation}
 S_{\mathrm{loop}}(V_{15})
 =S_{\mathrm{loop}}(V_q)+S_{\mathrm{loop}}(V_\ell)
 =\frac4{35}+\frac1{35}=\frac17.
 \label{eq:v6-component-boundary-action-additivity}
\end{equation}
Конечноранговый тепловой и детерминантный отклики раскладываются так же.
При $t=1$ получено
\begin{equation}
 0.07224234958+0.01806058740=0.09030293698,
\end{equation}
то есть ровно полный отклик ранга пятнадцать без перекрёстного члена.

В существующем действии нет пространственной координаты, измеряющей
расстояние между ядрами $q_q$ и $q_\ell$, и нет оператора, связывающего
кварковую и лептонную компоненты. Поэтому топология требует суммарный
класс, но не требует совместного центра. Два дефекта могут быть
алгебраически сложены, не образуя динамически связанную частицу.

\section{Компонентные щели}

Скалярный локальный перенос ограничивается на обе компоненты:
\begin{equation}
 W_s(k)=W_{\mathrm{dir}}(k)\otimes I_{H_s}.
\end{equation}
Если сохранять общий перенос, остаётся один свободный параметр смешения
$m$. Если разрешить независимые переносы, появляются два параметра
$m_q,m_\ell$. Ни унитарность, ни след не выбирают ненулевое значение.

Формальные подстановки
\begin{equation}
 m_q^2=\frac4{35},
 \qquad
 m_\ell^2=\frac1{35}
\end{equation}
совместимы с унитарностью, но не следуют из неё. Они были бы новым
правилом отождествления диагонального веса с внедиагональной амплитудой.
Следовательно, ни кварковая, ни лептонная массовая щель не выведена.

\begin{theorem}[Граница компонентного разложения]
Кварковая опора ранга двенадцать и лептонная опора ранга три являются
настоящими редуцирующими конечными, тёплицевыми, неограниченными и
Real-$K$-теоретическими подциклами. Однако ни один из них отдельно не
представляет исходный вынужденный класс пятнадцать. Их сумма восстанавливает
глобальный класс, но текущие следовые действия аддитивны и не создают ни
совместной локализации, ни энергии связывания, ни ненулевой массовой щели.
\end{theorem}

\begin{remark}
Тем самым фраза «пакет одного поколения» приобретает точный смысл: это
обязательная сумма классов в глобальном ledger, но пока не единая
локализованная частица. Нельзя отождествлять топологическую сумму с
физическим связанным состоянием.
\end{remark}

\section{Следующий различающий тест}

Следующий гейт о совместной локализации компонент должен проверить весь
уже существующий пространственный родитель на
перекрёстный член между $q_q(x)$ и $q_\ell(x)$:
\begin{enumerate}
 \item общую gauge- или хиггсовскую связность;
 \item смешанную кривизну моритова соответствия;
 \item общую моду порядка $Q,T,B$ или дефектный профиль;
 \item индексное условие рождения Real-пары в разных пространственных
 центрах.
\end{enumerate}
Стоп-критерий: если все существующие функционалы остаются прямой суммой,
совместная локализация классов $12$ и $3$ не выведена, и следующий
материальный объект следует искать внутри компонент, а не называть весь
класс пятнадцать частицей.

\begin{protocol}
\begin{enumerate}
 \item конечные редуцирующие подциклы рангов $12$ и $3$:
 \textbf{существуют};
 \item явные KO7-символы компонент: \textbf{построены};
 \item KO6-классы компонент: \textbf{$12$ и $3$};
 \item неограниченные тёплицевы циклы: \textbf{существуют};
 \item каждая компонента отдельно равна исходному классу $15$:
 \textbf{нет};
 \item прямая сумма восстанавливает класс $15$: \textbf{да};
 \item смешанный член петлевого действия: \textbf{ноль};
 \item топология принуждает совместный пространственный центр:
 \textbf{нет};
 \item ненулевая компонентная щель: \textbf{не выведена};
 \item класс $15$ является доказанной связанной частицей: \textbf{нет};
 \item следующий тест: \textbf{пространственная колокализация компонент}.
\end{enumerate}
\end{protocol}

Машинный сертификат сохранён в
\path{s2t_v6_spectral_transition_component_boundary_gate_results.json}.

\begin{thebibliography}{9}
\bibitem{component-boundary-krajewski1998}
T.~Krajewski,
\emph{Classification of Finite Spectral Triples},
Journal of Geometry and Physics 28 (1998), 1--30;
arXiv:hep-th/9701081.

\bibitem{component-boundary-atiyah1966}
M.~F.~Atiyah,
\emph{$K$-Theory and Reality},
Quarterly Journal of Mathematics 17 (1966), 367--386.

\bibitem{component-boundary-pimsner-voiculescu1980}
M.~Pimsner, D.~Voiculescu,
\emph{Exact Sequences for $K$-Groups and Ext-Groups of Certain
Cross-Product $C^*$-Algebras},
Journal of Operator Theory 4 (1980), 93--118.

\bibitem{component-boundary-baaj-julg1983}
S.~Baaj, P.~Julg,
\emph{Théorie bivariante de Kasparov et opérateurs non bornés dans les
modules hilbertiens},
Comptes Rendus de l'Académie des Sciences de Paris, Série I 296 (1983),
875--878.
\end{thebibliography}